\begin{statement}[C.8.3; Proved: the comparison of the two local discs]
Every complex number $t_0$ with $|t_0-1|\le1/500$ has a
representative $\tau$ with $|\tau-i|<1/50$ and $j(\tau)=1728t_0$.
Consequently every zero of $\wp$ for that modulus satisfies
$N\le\sqrt2$, with equality only for $t_0=1$.
\end{statement}
\begin{proof}
At $i$, $E_6=0$ and the Ramanujan differential identity gives
$|E_6'(i)|=\pi E_4(i)^2>13/2$, since $E_4(i)>29/20$.
On $|\delta|=1/50$, \eqref{eq:ca-derivatives} gives
\[
 |E_6(\tau)|>\frac{13}{100}-\frac1{50}=\frac{11}{100}.
\]
On the whole disc, $|q|<1/470$ and
$|\Delta_0|\le |q|\exp(24|q|/(1-|q|))<1/400$.
It follows on the boundary that
\[
 |t-1|=\frac{|E_6|^2}{1728|\Delta_0|}
 >\frac{121}{43200}>\frac1{500}.
\]
Rouch\'e applied to $E_6'(i)\delta$ shows that $E_6$ has
exactly one zero in the disc. Hence $t-1$ has two zeros there,
with multiplicity, and another application of Rouch\'e shows
that $t-t_0$ has two for each $|t_0-1|\le1/500$.
This argument proves coverage of the $t$-disc, rather than
inferring coverage merely from a boundary bound. Elliptic
curves over $\C$ with the same $j$ are isomorphic. In terms of
lattices the isomorphism is a homothety; it carries the two
zeros of $\wp$ to the two zeros of the scaled $\wp$ and
preserves $N$. Statement C.8.2 applies to either zero.
\end{proof}
